% Spellbook template for D&D 5e spells with level grouping
% Uses dndbook document class with structured spell organization

% Base LaTeX template for D&D 5e documents using DND-5e-LaTeX-Template
% This template provides the foundation for all D&D-style documents with structured content

\documentclass[letterpaper, 11pt, bg=full, twocolumn, stats=modern]{dndbook}

% DND-5e-LaTeX-Template provides most styling automatically
% Only include additional packages if needed

% Additional packages for enhanced functionality
\usepackage{graphicx}
\usepackage{wrapfig}
\usepackage{hyperref}
\usepackage{bookmark}
\usepackage{makeidx}

% Additional packages for spellbook layout
\usepackage{multicol}
\usepackage{enumitem}
\usepackage{tcolorbox}
\usepackage{titlesec}

% Enhanced list environments for spell details
\tcbuselibrary{breakable}

\title{Spell Collection}
\date{\today}

% Enable automatic numbering for chapters only, not sections
\setcounter{secnumdepth}{0}

% Disable text justification but preserve paragraph indentation
\raggedright
\raggedbottom
% Restore paragraph indentation after \raggedright resets it
\setlength{\parindent}{0.8em}

% Hyperlink configuration
\hypersetup{
    colorlinks=true,
    linkcolor=black,
    urlcolor=black,
    citecolor=black,
    filecolor=black,
    bookmarks=true,
    bookmarksopen=true,
    bookmarksnumbered=true,
    pdftitle={Spell Collection},
    pdfsubject={D&D 5e Content},
    pdfkeywords={D&D, 5e, RPG, tabletop, }
}

% Enhanced section styling for spellbook
% DND-5e-LaTeX-Template provides all styling automatically

% Configure for spellbook structure
\setcounter{secnumdepth}{0}  % No section numbers for cleaner look
\setcounter{tocdepth}{2}     % Include subsections in table of contents

% Studiorum Image Helpers — simple, explicit commands
% Load required packages locally so templates don't need to manage them.
\RequirePackage{graphicx}
\RequirePackage{xparse}
\RequirePackage{caption}

% Inline (non-floating) image block
% Usage: \StudiorumImageInline[<width fraction>]{file}{caption}[label]
\NewDocumentCommand{\StudiorumImageInline}{ s o m m o }{%%
  \par\medskip
  \noindent\centering
  \IfNoValueTF{#2}{%%
    \includegraphics[width=\linewidth]{#3}%%
  }{%%
    \includegraphics[width=#2\linewidth]{#3}%%
  }%%
  \IfBooleanF{#1}{%%
    \par\smallskip
    \captionof{figure}{#4}%%
    \IfNoValueF{#5}{\label{#5}}%%
  }%%
  \par\medskip
}

% Standard floating figure
% Usage: \StudiorumImageFloat[<width fraction>]{file}{caption}[label]
\NewDocumentCommand{\StudiorumImageFloat}{ s o m m o }{%%
  
\begin{figure}[htbp]
    \centering
    \IfNoValueTF{#2}{%%
      \includegraphics[width=\linewidth]{#3}%%
    }{%%
      \includegraphics[width=#2\linewidth]{#3}%%
    }%%
    \IfBooleanF{#1}{%%
      \caption{#4}%%
      \IfNoValueF{#5}{\label{#5}}%%
    }%%
  \end{figure}
%%
}

% Wide (two-column spanning when applicable)
% Usage: \StudiorumImageWide[<width fraction>]{file}{caption}[label]
\NewDocumentCommand{\StudiorumImageWide}{ s o m m o }{%%
  \ifdim\columnwidth=\textwidth
    
\begin{figure}[htbp]
      \centering
      \IfNoValueTF{#2}{%%
        \includegraphics[width=\linewidth]{#3}%%
      }{%%
        \includegraphics[width=#2\linewidth]{#3}%%
      }%%
      \IfBooleanF{#1}{%%
        \caption{#4}%%
        \IfNoValueF{#5}{\label{#5}}%%
      }%%
    \end{figure}
%%
  \else
    
\begin{figure*}[tp]
      \centering
      \IfNoValueTF{#2}{%%
        \includegraphics[width=\textwidth]{#3}%%
      }{%%
        \includegraphics[width=#2\textwidth]{#3}%%
      }%%
      \IfBooleanF{#1}{%%
        \caption{#4}%%
        \IfNoValueF{#5}{\label{#5}}%%
      }%%
    \end{figure*}
%%
  \fi
}

% Full-page (dedicated page; fits within text block)
% Usage: \StudiorumImageFullpage[<width fraction>]{file}{caption}[label]
\NewDocumentCommand{\StudiorumImageFullpage}{ s o m m o }{%%
  \clearpage
  
\begin{figure}[p]
    \centering
    \IfNoValueTF{#2}{%%
      \includegraphics[width=\textwidth,height=\textheight,keepaspectratio]{#3}%%
    }{%%
      \includegraphics[width=#2\textwidth,height=\textheight,keepaspectratio]{#3}%%
    }%%
    \IfBooleanF{#1}{%%
      \caption{#4}%%
      \IfNoValueF{#5}{\label{#5}}%%
    }%%
  \end{figure}

  \clearpage
}

% Full-page bleed (fills the physical page; may distort aspect ratio)
% Usage: \StudiorumImageFullpageBleed{file}{caption}[label]
\NewDocumentCommand{\StudiorumImageFullpageBleed}{ s o m m o }{%%
  \clearpage
  
\begin{figure}[p]
    \centering
    \includegraphics[width=\paperwidth,height=\paperheight]{#3}%%
    \IfBooleanF{#1}{%%
      \caption{#4}%%
      \IfNoValueF{#5}{\label{#5}}%%
    }%%
  \end{figure}

  \clearpage
}

% No-caption convenience aliases
\NewDocumentCommand{\StudiorumImageInlineNoCaption}{ o m }{\StudiorumImageInline*[#1]{#2}{}}
\NewDocumentCommand{\StudiorumImageFloatNoCaption}{ o m }{\StudiorumImageFloat*[#1]{#2}{}}
\NewDocumentCommand{\StudiorumImageWideNoCaption}{ o m }{\StudiorumImageWide*[#1]{#2}{}}
\NewDocumentCommand{\StudiorumImageFullpageNoCaption}{ o m }{\StudiorumImageFullpage*[#1]{#2}{}}
\NewDocumentCommand{\StudiorumImageFullpageBleedNoCaption}{ m }{\StudiorumImageFullpageBleed*{#1}{}}

\begin{document}

\tableofcontents
\clearpage

\section{1st Level Spells}

\vspace{0.5em}

\addcontentsline{toc}{subsection}{Shield}
\DndSpellHeader%
  {Shield}
  {\textit{1st-level abjuration}
}
  {1 reaction (which you take when you are hit by an attack or targeted by the \{@spell magic missile\} spell)}
  {Touch}
  {  V, S
}
  {1 round}

\noindent
An invisible barrier of magical force appears and protects you. Until the start of your next turn, you have a +5 bonus to AC, including against the triggering attack, and you take no damage from \textit{magic missile}.

\vspace{1.5em}

\section{3rd Level Spells}

\vspace{0.5em}

\addcontentsline{toc}{subsection}{Fireball}
\DndSpellHeader%
  {Fireball}
  {\textit{3rd-level evocation}
}
  {1 action}
  {150 feet}
  {  V, S, M (a tiny ball of bat guano and sulfur)
}
  {Instantaneous}

\noindent
A bright streak flashes from your pointing finger to a point you choose within range and then blossoms with a low roar into an explosion of flame. Each creature in a 20-foot-radius sphere centered on that point must make a Dexterity saving throw. A target takes 8d6 fire damage on a failed save, or half as much damage on a successful one.

The fire spreads around corners. It ignites flammable objects in the area that aren't being worn or carried.

  \paragraph{At Higher Levels}

When you cast this spell using a spell slot of 4th level or higher, the damage increases by 8d6 for each slot level above 3rd.

\vspace{1.5em}

\section{9th Level Spells}

\vspace{0.5em}

\addcontentsline{toc}{subsection}{Wish}
\DndSpellHeader%
  {Wish}
  {\textit{9th-level conjuration}
}
  {1 action}
  {Touch}
  {  V
}
  {Instantaneous}

\noindent
Wish is the mightiest spell a mortal creature can cast. By simply speaking aloud, you can alter the very foundations of reality in accord with your desires.

The basic use of this spell is to duplicate any other spell of 8th level or lower. You don't need to meet any requirements in that spell, including costly components. The spell simply takes effect. Alternatively, you can create one of the following effects of your choice:

\begin{itemize}
\item You create one object of up to 25,000 gp in value that isn't a magic item. The object can be no more than 300 feet in any dimension, and it appears in an unoccupied space you can see on the ground.
\item You allow up to twenty creatures that you can see to regain all hit points, and you end all effects on them described in the \textit{greater restoration} spell.
\item You grant up to ten creatures that you can see resistance to a damage type you choose.
\item You grant up to ten creatures you can see immunity to a single spell or other magical effect for 8 hours. For instance, you could make yourself and all your companions immune to a \textbf{lich}'s life drain attack.
\item You undo a single recent event by forcing a reroll of any roll made within the last round (including your last turn). Reality reshapes itself to accommodate the new result. For example, a \textit{wish} spell could undo an opponent's successful save, a foe's critical hit, or a friend's failed save. You can force the reroll to be made with advantage or disadvantage, and you can choose whether to use the reroll or the original roll.
\end{itemize}

You might be able to achieve something beyond the scope of the above examples. State your wish to the DM as precisely as possible. The DM has great latitude in ruling what occurs in such an instance; the greater the wish, the greater the likelihood that something goes wrong. This spell might simply fail, the effect you desire might only be partly achieved, or you might suffer some unforeseen consequence as a result of how you worded the wish. For example, wishing that a villain were dead might propel you forward in time to a period when that villain is no longer alive, effectively removing you from the game. Similarly, wishing for a legendary magic item or artifact might instantly transport you to the presence of the item's current owner.

The stress of casting this spell to produce any effect other than duplicating another spell weakens you. After enduring that stress, each time you cast a spell until you finish a long rest, you take 1d10 necrotic damage per level of that spell. This damage can't be reduced or prevented in any way. In addition, your Strength drops to 3, if it isn't 3 or lower already, for 2d4 days. For each of those days that you spend resting and doing nothing more than light activity, your remaining recovery time decreases by 2 days. Finally, there is a 33 percent chance that you are unable to cast wish ever again if you suffer this stress.

\end{document}